\documentclass[10pt]{article}
\usepackage[a4paper,margin=15mm]{geometry}
\usepackage[no-math]{fontspec}\setmainfont{Latin Modern Roman}
\usepackage{amsmath,amssymb,amsthm,mathtools,graphicx,xcolor,hyperref,xurl,enumitem}
\setlength{\emergencystretch}{4em}
\newtheorem{thm}{Theorem}[section]
\newtheorem{prop}[thm]{Proposition}
\usepackage{bm}
\usepackage[mathscr]{eucal}

\newcommand{\RR}{\mathbb{R}}
\newcommand{\ZZ}{\mathbb{Z}}
\newcommand{\NN}{\mathbb{N}}
\newcommand{\QQ}{\mathbb{Q}}
\newcommand{\CC}{\mathbb{C}}


\newcommand{\cA}{\mathcal{A}}
\newcommand{\cC}{\mathcal{C}}
\newcommand{\cF}{\mathcal{F}}
\newcommand{\cG}{\mathcal{G}}
\newcommand{\cP}{\mathcal{P}}
\newcommand{\cR}{\mathcal{R}}
\newcommand{\cU}{\la{U}}
\newcommand{\cY}{\mathcal{Y}}

\newcommand{\sC}{\mathscr{C}}
\newcommand{\sN}{\mathscr{N}}
\newcommand{\sO}{\mathscr{O}}
\newcommand{\sW}{\mathscr{W}}
\newcommand{\sR}{\mathscr{R}}
\newcommand{\sS}{\mathscr{S}}

\DeclareMathOperator{\ch}{ch}
\DeclareMathOperator{\wt}{wt}
\newcommand{\ii}{{\bf{i}}}

\allowdisplaybreaks

\usepackage{array}
\begin{document}
\section*{1259. Two-variable mod-ten quadruple q-series factor into independent Rogers--Ramanujan sums}
Katherine Baker, Shashank Kanade, Matthew C. Russell and Christopher Sadowski. \emph{Principal subspaces of basic modules for twisted affine Lie algebras, q-series multisums, and Nandi's identities}, Algebraic Combinatorics 6(6)(2023), 1533--1556. DOI \url{https://doi.org/10.5802/alco.311}. Exact archived official native publisher version, publication posted 2024-01-08. Native and original publisher PDF cover: CC BY 4.0, \url{https://creativecommons.org/licenses/by/4.0/}.
\paragraph{Selected problem (editorial).} Prove the exact formal three-variable factorization (96), with the complete recurrence and uniqueness theorem 6.2 and the factorization part of the proof of theorem 6.3. This is one outcome, with all four fundamental relations, the printed nine-term certificate, and coefficientwise induction kept as new core. The numerical specializations (79)--(82), (97)--(100) and representation-theoretic principal-character equalities are retained author context and receive no additional count. The selected statement is the author's exact displayed (96), not a stronger identity.
\paragraph{Ordinary prerequisites (editorial).} Formal power-series algebra and standard q-Pochhammer notation are allowed. The cited ordinary one-variable Rogers--Ramanujan series recurrence from Andrews, chapter 7, and the classical Rogers--Ramanujan product identities are ordinary prerequisites; the new four-variable recurrence/certificate and uniqueness are not assumed. The unselected Nandi/Murray--Miller algorithms and large external computer certificates are neither required nor included.
\paragraph{Difficulty (editorial).} H1. The selected formal factorization still requires constructing the nonroutine two-variable refinement and printed telescoping nine-term certificate for an exceptional four-variable quadratic sum. After these new equations are established, the coefficientwise uniqueness and one-variable Rogers--Ramanujan closing are routine, so the outcome is conservatively above ordinary qualification level. The certificate/refinement remains required core rather than a prerequisite; native 1063--1169, published physical 19--20.
\paragraph{Formatting (editorial).} Original author formulas, prose, comments, source order and theorem/equation indices are retained in the selected blocks. The original unused bbm import is omitted: no mathbbm token occurs in the complete article. No supplementary code is executed or redistributed. Standalone page flow and typography differ from the publisher. The original introduction's equation (8), referred to in the section introduction, is supplied verbatim as unselected source context below; equation (7) is retained alongside it only to preserve its actual shared align environment.
\clearpage
\section*{Original introductory context}
\setcounter{equation}{6}
\begin{align}
\sum_{i,j,k,\ell \ge 0}
\frac{q^{4i^2 + 12 ij + 8ik + 4i\ell + 12j^2 + 16jk + 8j\ell + 6k^2+6k\ell+2\ell^2}}
{\left(q^2;q^2\right)_i\left(q^2;q^2\right)_j\left(q;q\right)_k\left(q;q\right)_\ell}
&= \left(q^2,q^3,q^4,q^{10},q^{11},q^{12};q^{14}\right)^{-1}_\infty \label{eqn:Nandi_intro}, \\ 
\sum_{i,j,k,\ell \ge 0}
\frac{q^{2i^2 + 6 ij + 4ik + 2i\ell + 6j^2 + 8 jk + 4 j\ell + 3k^2+3k\ell+\ell^2}}
{\left(q^2;q^2\right)_i\left(q^2;q^2\right)_j\left(q;q\right)_k\left(q;q\right)_\ell} 
&=\left(q,q^2,q^4,q^6,q^8,q^9;q^{10}\right)^{-1}_\infty\label{eqn:mod10_intro}.
\end{align}
\section*{Original q-series notation}
\setcounter{section}{2}
\setcounter{subsection}{1}
\setcounter{equation}{10}
\subsection{\texorpdfstring{$q$-Series}{q-Series}}
We shall use standard notation regarding $q$-series. 
All of our series are formal, and issues of analytic convergence are disregarded.

For $n\in\ZZ_{\geq 0}\cup\{\infty\}$ we define:
\begin{align}
(a;q)_n=\prod_{0\leq t<n}(1-aq^{t}).
\end{align}
We will simply write $(a)_n$ when the base $q$ is understood.
To compress notation, we will write
\begin{align}
(a_1,a_2,\dots,a_r;q)_n=(a_1;q)_n(a_2;q)_n\cdots(a_r;q)_n.
\end{align}
We will also use modified theta functions:
\begin{align}
\theta(a;q)=(a;q)_\infty(q/a;q)_\infty,
\end{align}
and
\begin{align}
\theta(a_1,a_2,\dots,a_r;q)=\theta(a_1;q)\theta(a_2;q)\cdots(a_r;q).
\end{align}

\setcounter{section}{5}
\setcounter{equation}{78}
\section{New mod 10 identities}
\label{sec:mod10}

Varying the linear terms in the exponent of $q$ from~\eqref{eqn:mod10_intro}, we have the following list of identities where the products are related to the
principal characters standard (i.e., highest-weight, integrable) $\mathrm{D}_4^{(3)}$ modules of level $4$.
For notation regarding $\mathrm{D}_4^{(3)}$, see Carter's book \cite[P.\ 608]{Car-book}. 
For more information on principal characters and computational techniques for them, see the work of Bos~\cite{Bos}.
\begin{align}
&\sum_{i,j,k,\ell \ge 0}
\frac{q^{2i^2 + 6 ij + 4ik + 2i\ell + 6j^2 + 8 jk + 4 j\ell + 3k^2+3k\ell+\ell^2}}
{\left(q^2;q^2\right)_i\left(q^2;q^2\right)_j\left(q;q\right)_k\left(q;q\right)_\ell}\nonumber\\
&\quad=
\dfrac{1}{\theta(q,q^2,q^4;q^{10})} 
=\chi(\Omega(2\Lambda_0+\Lambda_1)),
\label{eq:mod10A}
\\
&\sum_{i,j,k,\ell \ge 0}
\frac{q^{2i^2 + 6 ij + 4ik + 2i\ell + 6j^2 + 8 jk + 4 j\ell + 3k^2+3k\ell+\ell^2+2j+k+\ell}}
{\left(q^2;q^2\right)_i\left(q^2;q^2\right)_j\left(q;q\right)_k\left(q;q\right)_\ell}
\nonumber\\ 
&\quad=
\dfrac{1}{\theta(q^2,q^2,q^3;q^{10})}  
=\chi(\Omega(2\Lambda_1)),
\\
&\sum_{i,j,k,\ell \ge 0}
\frac{q^{2i^2 + 6 ij + 4ik + 2i\ell + 6j^2 + 8 jk + 4 j\ell + 3k^2+3k\ell+\ell^2+2i+2j+2k}}
{\left(q^2;q^2\right)_i\left(q^2;q^2\right)_j\left(q;q\right)_k\left(q;q\right)_\ell}
\nonumber\\ 
&\quad= 
\dfrac{1}{\theta(q,q^4,q^4;q^{10})}  
=\chi(\Omega(\Lambda_0+\Lambda_2)),
\\
&\sum_{i,j,k,\ell \ge 0}
\frac{q^{2i^2 + 6 ij + 4ik + 2i\ell + 6j^2 + 8 jk + 4 j\ell + 3k^2+3k\ell+\ell^2+2i+4j+3k+\ell}}
{\left(q^2;q^2\right)_i\left(q^2;q^2\right)_j\left(q;q\right)_k\left(q;q\right)_\ell}
\nonumber\\ 
&\quad=
\dfrac{1}{\theta(q^2,q^3,q^4;q^{10})}   
= \chi(\Omega(4\Lambda_0)). \label{eq:mod10D}
\end{align}
Recall that  identities for principal characters of  standard modules
of $\mathrm{D}_4^{(3)}$ with level $3$ were previously conjectured by the second and third authors~\cite{KanRus-idf}, and
analytic forms of these conjectures were found by Kur\c{s}ung\"{o}z~\cite{Kur-KR}. See also the thesis of the third author~\cite{Rus-thesis} for further
related identities.

We will devote the remainder of this section to a proof of these conjectures.

As in the previous section, we begin by introducing refinements of the quadruple sums, along with arbitrary linear terms in the exponent of $q$. However, this time, we use two variables, $x$ and $y$.
\begin{align}
\label{eqn:defR}
&R_{A,B,C,D}(x,y,q)\\
&= \sum_{i,j,k,\ell \ge 0}
\frac{x ^{2j + k + \ell } y ^{ i+j+k}q^{2i^2 + 6 ij + 4ik + 2i\ell + 6j^2 + 8 jk + 4 j\ell + 3k^2+3k\ell+\ell^2+Ai+Bj+Ck+D\ell}}
{\left(q^2;q^2\right)_i\left(q^2;q^2\right)_j\left(q;q\right)_k\left(q;q\right)_\ell}. \nonumber
\end{align}
As usual, we will drop the arguments $x,y,q$ when they are clear.

We have the following shifts of $R_{A,B,C,D}$:
\begin{align}
R_{A,B,C,D}(xq,y,q)&=R_{A,B+2,C+1,D+1}(x,y,q) \label{eqn:shiftRx},\\
R_{A,B,C,D}(x,yq,q)&=R_{A+1,B+1,C+1,D}(x,y,q) \label{eqn:shiftRy}.
\end{align}


The fundamental relations governing these sums are deduced easily.
\begin{prop}
\label{prop:mod10rels}
\begin{alignat}{3}
m_1(A,B,C,D):&\quad R_{A,B,C,D}-R_{A+2,B,C,D}-yq^{2+A}R_{A+4,B+6,C+4,D+2}&&=0 \label{eqn:relm1},\\
m_2(A,B,C,D):&\quad R_{A,B,C,D}-R_{A,B+2,C,D}-x^2yq^{6+B}R_{A+6,B+12,C+8,D+4}&&=0 \label{eqn:relm2},\\
m_3(A,B,C,D):&\quad R_{A,B,C,D}-R_{A,B,C+1,D}-xyq^{3+C}R_{A+4,B+8,C+6,D+3}&&=0 \label{eqn:relm3},\\
m_4(A,B,C,D):&\quad R_{A,B,C,D}-R_{A,B,C,D+1}-xq^{1+D}R_{A+2,B+4,C+3,D+2}&&=0 \label{eqn:relm4}.
\end{alignat}
\end{prop}


\begin{thm}
The series $R_{0,0,0,0}$ is the unique solution in $\ZZ[[x,y,q]]$ satisfying
the following:
\begin{align}
F(x,y,q) &= F(xq,y,q) + xqF\left(xq^2,y,q\right) \label{eqn:m10d1},\\
F(x,y,0) &= 1,\quad 
F(x,0,q) = \sum_{\ell\geq 0}\dfrac{q^{\ell^2}x^\ell}{(q)_\ell},\quad 
F(0,y,q) = \sum_{i\geq 0}\dfrac{q^{2i^2}y^i}{(q^2;q^2)_i}. \label{eqn:m10init}
\end{align}
\end{thm}
\begin{proof}
$R_{0,0,0,0}$ is seen to satisfy \eqref{eqn:m10init} easily.
Due to the shifts \eqref{eqn:shiftRx} and \eqref{eqn:shiftRy}, showing that $R_{0,0,0,0}$ satisfies
\eqref{eqn:m10d1} is equivalent to showing:
\begin{align}
R_{0,0,0,0} - R_{0,2,1,1} - xqR_{0,4,2,2} = 0.
\end{align}
This relation can be obtained as a linear combination of the fundamental relations \eqref{eqn:relm1}--\eqref{eqn:relm4}:
\begin{align}
-&m_1(-2,0,0,0)+m_1(-2,0,0,1)-xq\cdot m_1(0,4,2,2)
\nonumber\\
&+xq\cdot m_1(0,4,3,2)+m_2(0,0,0,1)+m_3(0,2,0,1)
\nonumber \\
&-xq\cdot m_3(2,4,2,2)+m_4(-2,0,0,0)-y\cdot m_4(2,6,4,2).
\end{align}

Now we prove the uniqueness. 
Suppose that $F(x,y,q)=\sum_{i,j,k\geq 0}f_{i,j,k}x^iy^jq^k$ is a solution to this system.
Note that for all $i,j,k\geq 0$, $f_{0,j,k}$, $f_{i,0,k}$ and $f_{i,j,0}$ 
are uniquely determined due to the initial conditions 
\eqref{eqn:m10init}.
For convenience, we define $f_{i,j,k}=0$ whenever any one or more of $i,j,k$ are negative.
\eqref{eqn:m10d1} translates to:
\begin{align}
f_{i,j,k}&=f_{i,j,k-i}+f_{i-1,j,k-2i+1}.
\label{eqn:frel}
\end{align}
Now we induct on $N=i+k$ and show that all $f_{i,j,k}$ are uniquely determined.
The case~$N=0$ follows easily since $f_{0,j,0}$ are uniquely known due to \eqref{eqn:m10init}.
Suppose that for all triples $(i,j,k)$ with $i+k<N$, the values $f_{i,j,k}$ are determined.
Pick a triple~$(I,J,K)$ such that $I+K=N$. 
If $I=0$, we know $f_{0,J,N}$ by \eqref{eqn:m10init}.
So, suppose that~$I\geq 1$. 
Then, the RHS of \eqref{eqn:frel} involves two terms:
\begin{enumerate}
\item $f_{I,J,K-I}$ for which $I+(K-I) = K < N$, (since $K+I=N$ and $I\geq 1$)
\item $f_{I-1,J,K-2I+1}$ for which $(I-1)+(K-2I+1) = K - I < N$ (since $K+I=N$ and $I\geq 1$).
\end{enumerate}
This means that the RHS of \eqref{eqn:frel} has been determined already, and this determines
the LHS uniquely.
\end{proof}

\begin{thm}
Identities \eqref{eq:mod10A}--\eqref{eq:mod10D} are true.
\end{thm}
\begin{proof}
It is easy to see using chapter 7 of Andrews' text~\cite{And-book} that 
the series  
\begin{align}
\left(\sum_{i\geq 0} \dfrac{x^iq^{i^2}}{(q)_i}\right)
\left(\sum_{j\geq 0} \dfrac{y^jq^{2j^2}}{(q^2;q^2)_j}\right)
\end{align}
satisfies \eqref{eqn:m10d1} and \eqref{eqn:m10init}.
Thus, due to uniqueness, we have:
\begin{align}
R_{0,0,0,0}(x,y,q)=\left(\sum_{i\geq 0} \dfrac{x^iq^{i^2}}{(q)_i}\right)
\left(\sum_{j\geq 0} \dfrac{y^jq^{2j^2}}{(q^2;q^2)_j}\right).
\end{align}
By the Rogers--Ramanujan identities \cite[Ch.\ 7]{And-book}, we now have:
\begin{align}
R_{0,0,0,0}(1,1,q)=\dfrac{1}{\theta(q;q^5)\theta(q^2;q^{10})},\\
R_{0,0,0,0}(q,1,q)=\dfrac{1}{\theta(q^2;q^5)\theta(q^2;q^{10})},\\
R_{0,0,0,0}(1,q^2,q)=\dfrac{1}{\theta(q;q^5)\theta(q^4;q^{10})},\\
R_{0,0,0,0}(q,q^2,q)=\dfrac{1}{\theta(q^2;q^5)\theta(q^4;q^{10})}.
\end{align}
This immediately proves~\eqref{eq:mod10A}--\eqref{eq:mod10D}.
\end{proof}
\begin{thebibliography}{99}
\bibitem{A}Andrews, George E. An analytic generalization of the Rogers–Ramanujan identities for odd moduli, Proc. Nat. Acad. Sci. U.S.A., Volume 71 (1974), pp. 4082-4085
\bibitem{And-mult}Andrews, George E. Multiple series Rogers–Ramanujan type identities, Pacific J. Math., Volume 114 (1984) no. 2, pp. 267-283
\bibitem{And-book}Andrews, George E. The theory of partitions, Cambridge Mathematical Library, Cambridge University Press, Cambridge, 1998, xvi+255 pages (Reprint of the 1976 original)
\bibitem{ASW}Andrews, George E.; Schilling, Anne; Warnaar, S. Ole An A 2  Bailey lemma and Rogers–Ramanujan-type identities, J. Amer. Math. Soc., Volume 12 (1999) no. 3, pp. 677-702
\bibitem{Bos}Bos, M. K. Coding the principal character formula for affine Kac–Moody Lie algebras, Math. Comp., Volume 72 (2003) no. 244, pp. 2001-2012
\bibitem{BIS}Bressoud, D.; Ismail, M. E. H.; Stanton, D. Change of base in Bailey pairs, Ramanujan J., Volume 4 (2000) no. 4, pp. 435-453
\bibitem{Br}Bressoud, David M. Analytic and combinatorial generalizations of the Rogers–Ramanujan identities, Mem. Amer. Math. Soc., Volume 24 (1980) no. 227, p. 54
\bibitem{BriJenMah}Bringmann, Kathrin; Jennings-Shaffer, Chris; Mahlburg, Karl Proofs and reductions of various conjectured partition identities of Kanade and Russell, J. Reine Angew. Math., Volume 766 (2020), pp. 109-135
\bibitem{CalLM1}Calinescu, C.; Lepowsky, J.; Milas, A. Vertex-algebraic structure of the principal subspaces of certain A 1 (1) -modules. I. Level one case, Internat. J. Math., Volume 19 (2008) no. 1, pp. 71-92
\bibitem{CalLM2}Calinescu, C.; Lepowsky, J.; Milas, A. Vertex-algebraic structure of the principal subspaces of certain A 1 (1) -modules. II. Higher-level case, J. Pure Appl. Algebra, Volume 212 (2008) no. 8, pp. 1928-1950
\bibitem{CalLM3}Calinescu, C.; Lepowsky, J.; Milas, A. Vertex-algebraic structure of the principal subspaces of level one modules for the untwisted affine Lie algebras of types A,D,E, J. Algebra, Volume 323 (2010) no. 1, pp. 167-192
\bibitem{CalLM4}Calinescu, C.; Lepowsky, J.; Milas, A. Vertex-algebraic structure of principal subspaces of standard A 2 (2) -modules, I, Internat. J. Math., Volume 25 (2014) no. 7, Paper no. 1450063, 44 pages
\bibitem{CalMPe}Calinescu, Corina; Milas, Antun; Penn, Michael Vertex algebraic structure of principal subspaces of basic A 2n (2) -modules, J. Pure Appl. Algebra, Volume 220 (2016) no. 5, pp. 1752-1784
\bibitem{CPS}Calinescu, Corina; Penn, Michael; Sadowski, Christopher Presentations of principal subspaces of higher level standard A 2 (2) -modules, Algebr. Represent. Theory, Volume 22 (2019) no. 6, pp. 1457-1478
\bibitem{CLM1}Capparelli, S.; Lepowsky, J.; Milas, A. The Rogers–Ramanujan recursion and intertwining operators, Commun. Contemp. Math., Volume 5 (2003) no. 6, pp. 947-966
\bibitem{CLM2}Capparelli, S.; Lepowsky, J.; Milas, A. The Rogers–Selberg recursions, the Gordon–Andrews identities and intertwining operators, Ramanujan J., Volume 12 (2006) no. 3, pp. 379-397
\bibitem{Cap-1}Capparelli, Stefano On some representations of twisted affine Lie algebras and combinatorial identities, J. Algebra, Volume 154 (1993) no. 2, pp. 335-355
\bibitem{Car-book}Carter, R. W. Lie algebras of finite and affine type, Cambridge Studies in Advanced Mathematics, 96, Cambridge University Press, Cambridge, 2005, xviii+632 pages
\bibitem{Che-linked}Chern, Shane Linked partition ideals, directed graphs and q-multi-summations, Electron. J. Combin., Volume 27 (2020) no. 3, Paper no. 3.33, 29 pages
\bibitem{CorDouUnc}Corteel, Sylvie; Dousse, Jehanne; Uncu, Ali Kemal Cylindric partitions and some new A 2  Rogers–Ramanujan identities, Proc. Amer. Math. Soc., Volume 150 (2022) no. 2, pp. 481-497
\bibitem{FS1}Feigin, Boris; Stoyanovsky, A. V. Quasi-particles models for the representations of Lie algebras and geometry of flag manifold, 1993
\bibitem{Ga}Garvan, Frank A q-product tutorial for a q-series MAPLE package, Sém. Lothar. Combin., Volume 42 (1999), Paper no. B42d, 27 pages The Andrews Festschrift (Maratea, 1998)
\bibitem{GriOnoWar}Griffin, Michael J.; Ono, Ken; Warnaar, S. Ole A framework of Rogers–Ramanujan identities and their arithmetic properties, Duke Math. J., Volume 165 (2016) no. 8, pp. 1475-1527
\bibitem{KanRus-idf}Kanade, Shashank; Russell, Matthew C. IdentityFinder and some new identities of Rogers–Ramanujan type, Exp. Math., Volume 24 (2015) no. 4, pp. 419-423
\bibitem{KanRus-stair}Kanade, Shashank; Russell, Matthew C. Staircases to analytic sum-sides for many new integer partition identities of Rogers–Ramanujan type, Electron. J. Combin., Volume 26 (2019) no. 1, Paper no. 1.6, 33 pages
\bibitem{KanRus-a22}Kanade, Shashank; Russell, Matthew C. On q-series for principal characters of standard A 2 (2) -modules, Adv. Math., Volume 400 (2022), Paper no. 108282, 24 pages
\bibitem{KanRus-cylindric}Kanade, Shashank; Russell, Matthew C. Completing the A2 Andrews–Schilling–Warnaar Identities, Int. Math. Res. Not. IMRN (2023) no. 20, pp. 17100-17155
\bibitem{Kon}Konenkov, Stepan Further q-reflections on the modulo 9 Kanade–Russell (conjectural) identities, 2022
\bibitem{Kur-Cap}Kurşungöz, Kağan Andrews–Gordon type series for Capparelli’s and Göllnitz–Gordon identities, J. Combin. Theory Ser. A, Volume 165 (2019), pp. 117-138
\bibitem{Kur-KR}Kurşungöz, Kağan Andrews–Gordon type series for Kanade–Russell conjectures, Ann. Comb., Volume 23 (2019) no. 3-4, pp. 835-888
\bibitem{L}Lepowsky, J. Calculus of twisted vertex operators, Proc. Nat. Acad. Sci. U.S.A., Volume 82 (1985) no. 24, pp. 8295-8299
\bibitem{LL}Lepowsky, James; Li, Haisheng Introduction to vertex operator algebras and their representations, Progress in Mathematics, 227, Birkhäuser Boston, Inc., Boston, MA, 2004, xiv+318 pages
\bibitem{Nan-thesis}Nandi, Debajyoti Partition identities arising from the standard A2 (2)-modules of level 4, ProQuest LLC, Ann Arbor, MI, 2014, 203 pages Thesis (Ph.D.)–Rutgers The State University of New Jersey - New Brunswick
\bibitem{Pau}Paule, Peter On identities of the Rogers–Ramanujan type, J. Math. Anal. Appl., Volume 107 (1985) no. 1, pp. 255-284
\bibitem{PS1}Penn, Michael; Sadowski, Christopher Vertex-algebraic structure of principal subspaces of basic D 4 (3) -modules, Ramanujan J., Volume 43 (2017) no. 3, pp. 571-617
\bibitem{PS2}Penn, Michael; Sadowski, Christopher Vertex-algebraic structure of principal subspaces of the basic modules for twisted affine Lie algebras of type A 2n-1 (2) , D n (2) , E 6 (2) , J. Algebra, Volume 496 (2018), pp. 242-291
\bibitem{PSW}Penn, Michael; Sadowski, Christopher; Webb, Gautam Principal subspaces of twisted modules for certain lattice vertex operator algebras, Internat. J. Math., Volume 30 (2019) no. 10, Paper no. 1950048, 47 pages
\bibitem{Ros}Rosengren, Hjalmar Proofs of some partition identities conjectured by Kanade and Russell, Ramanujan J., Volume 61 (2023) no. 1, pp. 295-317
\bibitem{Rus-thesis}Russell, Matthew Christopher Using experimental mathematics to conjecture and prove theorems in the theory of partitions and commutative and non-commutative recurrences, ProQuest LLC, Ann Arbor, MI, 2016, 74 pages Thesis (Ph.D.)–Rutgers The State University of New Jersey - New Brunswick
\bibitem{Sil-book}Sills, Andrew V. An invitation to the Rogers–Ramanujan identities, CRC Press, Boca Raton, FL, 2018, xx+233 pages (With a foreword by George E. Andrews)
\bibitem{S}Stembridge, John R. Hall–Littlewood functions, plane partitions, and the Rogers–Ramanujan identities, Trans. Amer. Math. Soc., Volume 319 (1990) no. 2, pp. 469-498
\bibitem{FS2}Stoyanovskiĭ, A. V.; Feĭgin, B. L. Functional models of the representations of current algebras, and semi-infinite Schubert cells, Funktsional. Anal. i Prilozhen., Volume 28 (1994) no. 1, p. 68-90, 96
\bibitem{Ta}Takenaka, Ryo Vertex algebraic construction of modules for twisted affine Lie algebras of type A 2l (2) , J. Pure Appl. Algebra, Volume 227 (2023) no. 4, Paper no. 107263, 33 pages
\bibitem{TakTsu-nandi}Takigiku, Motoki; Tsuchioka, Shunsuke A proof of conjectured partition identities of Nandi, 2020 (Forthcoming, Amer. J. Math)
\bibitem{TakTsu-a22}Takigiku, Motoki; Tsuchioka, Shunsuke Andrews–Gordon type series for the level 5 and 7 standard modules of the affine Lie algebra A 2 (2) , Proc. Amer. Math. Soc., Volume 149 (2021) no. 7, pp. 2763-2776
\bibitem{Tsu-a21level3}Tsuchioka, Shunsuke An example of A 2  Rogers–Ramanujan bipartition identities of level 3, 2022
\bibitem{Tsu-d43level3}Tsuchioka, Shunsuke A vertex operator reformulation of the Kanade–Russell conjecture modulo 9, 2022
\bibitem{UncZud}Uncu, Ali; Zudilin, Wadim Reflecting (on) the modulo 9 Kanade–Russell (conjectural) identities, Sém. Lothar. Combin., Volume 85 ([2020–2021]), Paper no. B85e, 17 pages
\bibitem{War1}Warnaar, S. Ole The generalized Borwein conjecture. II. Refined q-trinomial coefficients, Discrete Math., Volume 272 (2003) no. 2-3, pp. 215-258
\bibitem{War2}Warnaar, S. Ole The A 2  Andrews-Gordon identities and cylindric partitions, Trans. Amer. Math. Soc. Ser. B, Volume 10 (2023), pp. 715-765
\end{thebibliography}
\end{document}
